\scanpage{46}

% 1–3, heading and opening
\uline{Ex 3} Les \unsure{triples|triplets|théorie} de Dold-Puppe nous donnent une \deleted{\unsure{foncteur} équivalence de catégories} équivalence de catégories (pour tte catégorie \deleted{$\unsure{Ab_k|Ab_k^\bullet}\quad Ab_k$.} additive \unsure{avec} facteurs directs)

% 1–3, upper left margin
% slanted
\margin{\uline{NB} Le foncteur $DP$ est défini de \unsure{ple.|faç.} \unsure{fid.} pour des $B$ additives (\unsure{qui ont|karoubiennes} des facteurs directs)}

% 4, abbreviations in left margin
\margin{$ND=$ non dégénéré \quad $DP=$ Dold-Puppe \deleted{\unsure{de $Ab$}}}

% 4–5, equivalences
\diagram{\begin{tikzcd}[ampersand replacement=\&,column sep=small]
 B_{\bullet}\arrow[r,shift right=1.1ex,"{DP^{*}}"']\& B_{*}\arrow[l,bend right=35,"{ND_{\bullet}}"']\arrow[l,phantom,"\simeq"] \& \text{donc aussi}\ \substack{\text{(en appliquant}\\\text{à } B^\circ\text{)}}\&
 B^{\bullet}\arrow[r,shift right=1.1ex,"{DP_{*}}"']\& B^{*}\arrow[l,bend right=35,"{ND^{\bullet}}"']\arrow[l,phantom,"\simeq"]
\end{tikzcd}}

% 4–5, left margin under abbreviations
% slanted
\margin{équiv. compatibles aux foncteurs additifs}

% 6
en particulier

% 7
\[
 Ab_{k\,\bullet}\overset{\sim}{\longrightarrow}Ab_{k\,*},\qquad Ab_k^{\bullet}\overset{\sim}{\longrightarrow}Ab_k^{*}
\]

% 8–13
\deleted{Il est immédiat} \deleted{\unsure{Donc}.} l'image \unsure{essentielles|essentielle} de $\mathbb K_{\bullet}^{\deleted{\unsure{k}}}$ et $\mathbb K^{\bullet}_{\deleted{\unsure{k}}}$ par ces deux foncteurs sont notés $\mathbb K_{*}$ et $\mathbb K^{*}$. On a donc nos diagrammes d'équivalences canoniques

% 9–13, left margin
% slanted
\margin{$k$ est sous-entendu dans ces notations, sinon, c'est $\mathbb K_{!}(k)$ avec indices ou exposant $\bullet$ ou $*$.}

% 14–15, cancelled diagram
\diagram{\begin{tikzcd}[ampersand replacement=\&,execute at end picture={\draw (dlla.south west) -- (dula.north east) -- (dlla.south east) -- (dura.north west) -- (dlra.south west) -- (dura.east);}]
 |[alias=dula]|\deleted{Ab_{k\,\bullet}}\arrow[r,shift right,"{\deleted{DP}}"']\&|[alias=dura]|\deleted{Ab_{k\,*}}\arrow[l,shift right,"{\deleted{ND}}"']\\
 |[alias=dlla]|\deleted{Ab_k^{\bullet}}\arrow[r,shift right]\&|[alias=dlra]|\deleted{Ab_k^{*}}\arrow[l,shift right]
\end{tikzcd}}

% 16–20, five-object diagram
\diagram{\begin{tikzcd}[ampersand replacement=\&,row sep=2.4em,column sep=2.4em]
 \mathbb K_{\bullet}\arrow[rr,shift right,"{DP^{*}}"']\arrow[dd,leftrightarrow,"{\vee\,\circ}", "{\deleted{\unsure{\mathrm{can}}}}"'] \& \& \mathbb K_{*}\arrow[ll,shift right,"{ND_{\bullet}}"']\arrow[dd,leftrightarrow,"{\circ\,\vee}", "{\deleted{\unsure{\mathrm{can}}}}"']\\
 \&\mathbb K\arrow[ul,"{\alpha_{\bullet}}"]\arrow[ur,"{\alpha_{*}}"']\arrow[dl,"{\alpha^{\bullet}\ \circ}"]\arrow[dr,"{\alpha^{*}\ \circ}"']\&\\
 \mathbb K^{\bullet}\arrow[rr,shift right,"{DP_{*}}"']\&\&\mathbb K^{*}\arrow[ll,shift right,"{ND^{\bullet}}"']
\end{tikzcd}}

% 16–23, long marginal convention, left
\margin{Convention d'écriture. pour un objet $C$ de $K$, on note $C_{\bullet}, C_{*}, C^{\bullet}, C^{*}$ les objets correspondants obtenus en appliquant $\alpha_{\bullet},\alpha_{*},\alpha^{\bullet},\alpha^{*}$, et ses comp. sont notés $C_n,C_{[n]},C^n,C^{[n]}$ (Attention, \uline{ne pas} oublier le \uline{signe} $[\ ]$ et $\mathbb E_{\bullet}^{i},\mathbb E_{*}$ $\mathbb E^{\bullet},\mathbb E^{*}$, $\mathbb E^{i}$}

% 21–22
Les flèches \deleted{\unsure{sont}} marquées d'un signe $\circ$ sont des foncteurs \uline{contravariant}.

% 23–25
On a

\uline{Proposition 4} \deleted{\illegible{2}} On a $\forall i\geq0$ un isom canonique

% archive, bottom left
% archive sheet number: 45
